\documentclass[11pt,a4paper]{article}

\usepackage[T1]{fontenc}
\usepackage[utf8]{inputenc}
\usepackage{lmodern}
\usepackage[margin=2.4cm]{geometry}
\usepackage{microtype}

% Code fragments are rare and short. Where one appears, the code IS the
% argument, so it gets a slightly indented block and nothing else.
\newenvironment{code}%
  {\par\vspace{0.5em}\begingroup\footnotesize\ttfamily\leftskip=1.2em%
   \parskip=0pt\baselineskip=1.05\baselineskip\obeylines}%
  {\endgroup\par\vspace{0.5em}}

\setlength{\parskip}{0.55em}
\setlength{\parindent}{0pt}

\usepackage{titlesec}
\titleformat{\section}{\large\bfseries}{\thesection.}{0.6em}{}
\titlespacing*{\section}{0pt}{1.4em}{0.5em}
\titleformat{\subsection}{\normalsize\bfseries}{}{0em}{}
\titlespacing*{\subsection}{0pt}{0.9em}{0.3em}

\begin{document}

\begin{flushleft}
{\LARGE\bfseries Design rationale}\\[0.35em]
{\large KYC onboarding and compliance review desk}\\[0.6em]
{\normalsize Denis Xhabrahimi \quad·\quad September 2026}
\end{flushleft}
\rule{\textwidth}{0.4pt}

\vspace{0.6em}

This is a record of why the system is built the way it is, including the
options I rejected. It assumes you have not clicked the demo.

The system onboards a customer: it takes their details, sends them through an
identity-verification vendor, screens the name against sanctions lists, scores
the result, and either decides automatically or hands the case to a compliance
officer at a review desk. Two services --- a Next.js web application and a
Python worker --- share one PostgreSQL database and never call each other
directly. The interesting problem is not the happy path. It is that the vendor
drops webhooks, delivers them twice, and delivers them out of order, and the
system has to stay correct anyway, and be able to prove afterwards that it was.

Everything below is a decision I would have to defend.

\section{The principle underneath: fail closed}

I applied one rule four times before I noticed I was doing it, so it is worth
naming before the individual decisions: \textbf{prefer a visible failure to a
silent degradation}.

A sanctions screen that runs against half a list returns fewer matches and
looks like a clean result. A disclosure filter that does not recognise an event
type and shows it anyway leaks. A shared secret that is unset and therefore
lets everyone through is worse than one that locks everyone out. In each case
the failure that does not announce itself is the expensive one, because nobody
goes looking for it.

So: the sanctions loader stamps ``loaded'' \emph{last}, and the matcher refuses
a snapshot without that stamp, so a half-finished load is unusable rather than
quietly thin. The applicant-facing audit view is a whitelist that renders
nothing for an event type it does not know. The worker's HTTP endpoint refuses
every request when its secret is unconfigured. The audit table rejects UPDATE
and DELETE at the database level rather than by convention. None of those
choices makes the system more capable. Each makes its failures louder.

\section{The job queue is a Postgres table}

This is the decision I would defend hardest, because it looks like the naive
one.

\subsection{The problem}

When a verification result arrives by webhook, I have to do two things: record
that it arrived, and schedule the work it implies. If those two are not atomic,
there is no ordering of them that is correct. Enqueue first and the database
write can fail, leaving a job referring to a row that does not exist. Write
first and the enqueue can fail, leaving a result nobody will ever process ---
an applicant waiting forever with no error anywhere.

\subsection{What I chose, and what I rejected}

The queue is an ordinary table in the same database as the business data, so
enqueueing is part of the same transaction as the write that caused it. Either
both happen or neither does. There is no window.

I rejected Redis and Celery for exactly this. Both are better queues in
isolation --- faster, with richer scheduling and visibility. Both are a
\emph{second} system, which means a distributed transaction I would have to
fake with an outbox pattern: write to Postgres, write a row saying ``I still
need to tell Redis'', have something drain that row into Redis. At which point
I have built a database-backed queue anyway, plus Redis, and the database-backed
queue is the one carrying correctness. I would rather have the honest version
of that than the version with a dependency in front of it.

I rejected \texttt{pg-boss} for a different reason: it would have owned the
schema. The database is shared by TypeScript and Python, and I did not want a
JavaScript library's migrations deciding the shape of a table the Python worker
also writes to. Raw SQL in both services, with numbered migrations belonging to
neither, keeps that boundary clean.

\subsection{When this stops being a good idea}

At high throughput. Every claim is a write and a row-level lock; the table
accumulates dead tuples and leans on autovacuum. Somewhere in the low thousands
of jobs per second this becomes the bottleneck and a purpose-built broker
becomes worth its operational cost. This system processes a few hundred
applications a day. I am nowhere near that, and the honest version of this
argument includes the point at which I would change my mind.

\section{FOR UPDATE SKIP LOCKED}

With several workers running, two of them must never claim the same job. The
claim is one statement:

\begin{code}
with claimed as (
~~~~select id from jobs
~~~~~where status = 'queued' and run\_after <= now()
~~~~~order by run\_after, id
~~~~~for update skip locked
~~~~~limit 1
)
update jobs j
~~~set status = 'running', attempts = j.attempts + 1,
~~~~~~~locked\_at = now(), locked\_by = \%s
~~from claimed c where j.id = c.id
returning j.id, j.job\_type, j.payload, j.attempts, j.max\_attempts
\end{code}

\texttt{FOR UPDATE} takes a row lock. Without \texttt{SKIP LOCKED}, a second
worker reaching the same row \emph{blocks} until the first commits --- every
worker queued behind one row, which is the opposite of what a worker pool is
for. \texttt{SKIP LOCKED} changes one thing: a locked row is treated as though
it were not there, so the second worker takes the next one instead.

I did not want to assert this, so I measured the counterfactual. There is a
deliberately broken claim function in the codebase --- select, then update, with
no lock --- kept so the comparison is reproducible rather than rhetorical. Two
workers, 200 jobs, no lock: \textbf{380 executions of 200 jobs, 180 of them run
twice}. Re-running it today gives 196 to 199 duplicates of 200; the exact number
moves with machine speed, and that it moves is the point --- it is a race, not a
constant. With the lock, across many runs, the duplicate count is zero.

There is a test that asserts the broken version is \emph{still broken}. It looks
absurd until you consider how it would fail: someone reads that function, sees a
race condition, fixes it, and the comparison quietly becomes a claim about
nothing.

\subsection{Attempts increments at claim time, not on failure}

A job that crashes the process hard enough --- out of memory, a killed container
--- never reaches a failure handler, so a counter incremented on failure never
moves. That job would be retried forever, killing a worker each time. Counting
at claim time means every attempt is paid for whether or not anyone survived to
report it. A poison pill exhausts its attempts and is parked instead of
circulating.

\section{The lifecycle lives in a database trigger}

An application moves \texttt{started} → \texttt{submitted} → \texttt{checking} →
\texttt{screening} → \texttt{decided}. The rule that matters is that nothing
reaches \texttt{decided} except from \texttt{screening}: no application is ever
decided without having been screened.

I had that rule in Python. It was correct, and it was in the wrong place. The
same table is written by the TypeScript service and the Python worker, so a rule
held in one of them is a rule the other can break --- and the guarantee I
actually wanted to make was about the \emph{data}, not about one codebase's
discipline.

So it is a \texttt{BEFORE UPDATE} trigger raising \texttt{restrict\_violation}:

\begin{code}
if new.status is not distinct from old.status then
~~~~return new;  -- not every update is a transition
end if;

if not exists (select 1 from application\_transitions t
~~~~~~~~~~~~~~~where t.from\_status = old.status
~~~~~~~~~~~~~~~~~~and t.to\_status = new.status) then
~~~~raise exception 'illegal transition: \% -> \%', old.status, new.status
~~~~~~using errcode = 'restrict\_violation';
end if;
\end{code}

The permitted transitions are rows in a table rather than a hard-coded list, so
the rule is inspectable with a query and changing it is a migration --- a
reviewable commit --- rather than an edit to a constant.

This is the same argument as the one for migrations belonging to neither
service, which was a stated principle of the project from the start. I shipped
the original version not following it. The Python table is still there as a
mirror, with a comment saying the database is authoritative, because the worker
benefits from knowing a transition will be refused before attempting it.

\textbf{The cost:} the rule now lives somewhere a developer reading the
application code will not see it. That is a real downside and I think it is the
right trade, because the alternative distributes a data guarantee across two
languages and trusts both.

\section{The audit log is append-only, enforced by the database}

Every state change writes a row to \texttt{audit\_events}. That table rejects
UPDATE, DELETE and TRUNCATE through triggers.

``Append-only by convention'' is not append-only. It is a code review standard,
and the thing an auditor is worried about is precisely the case where someone
decided the convention did not apply. Enforcing it in the database means the
guarantee survives a developer with a good reason and a deadline.

It also binds me. The test suite cannot clear the audit table between runs, so
tests that count audit rows have to record a high-water mark first and count
only above it --- and a demo script that resets the queue cannot reset the
history. That friction is the guarantee being real rather than decorative.

\textbf{What it does not give me:} this is tamper-\emph{evident} within the
system, not tamper-proof. Anyone with superuser access can disable a trigger.
Real immutability means WORM storage or an append-only log outside the database,
and a hash chain linking each row to its predecessor would make undetected
mid-history edits infeasible. I have not built either. I would rather say that
than let ``append-only'' do more work in a conversation than it can carry.

\section{Risk scoring is points, not machine learning}

Each signal contributes a fixed number of points: a confirmed sanctions match
60, a politically-exposed-person match 30, residence in a jurisdiction subject
to a call for action 40, a failed document check 40. Under 20 auto-approves,
20--79 goes to a human, 80 and over auto-rejects.

The argument is not that a model would be less accurate. It is that when a
regulator, an ombudsman, or the customer asks why this person was refused, the
answer has to be a sentence. ``Sanctions match at 100\% name similarity,
downgraded to probable because the listed date of birth conflicts, plus a
declined document check'' is a sentence. ``The model scored 0.83'' is not, and
neither is a SHAP plot.

Two consequences I chose deliberately. PEP matches cannot on their own reach the
rejection threshold, because being a politically exposed person is not illegal
and wholesale refusal --- de-risking --- is something regulators criticise; the
law asks for enhanced due diligence, which means a human. And a sanctions match
alone routes to review rather than rejection, because the screening finds
\emph{names}, and identity is a judgement a person makes.

Every scored application stores the ruleset version and the thresholds in force
at the time, so a decision can be reconstructed after the rules have moved.

\textbf{The cost:} the weights are my judgement, not fitted to outcomes. A real
firm would back-test them and have a model risk governance process. A points
system is explainable and unvalidated, which is a genuinely different thing from
being right.

\section{Two audiences, two rules: the disclosure boundary}

The applicant's status page and the officer's case view read the same audit
table. They must not show the same thing.

The officer's view renders the raw event payload, because an auditor needs the
unvarnished record and a pretty-printer that dropped a field would defeat the
purpose of keeping one. The applicant's page is public --- addressed by the
application's UUID, no login --- and the same rows contain the risk score, the
routing decision, and sentences of the form \emph{``Sanctions match against X on
OFAC SDN at 86\% name similarity''}.

I originally shipped the applicant page rendering those rows directly. That was
the most serious mistake in the project.

Three reasons an applicant is not shown that, none of which needs a citation to
stand up. First, tipping-off: financial-crime regimes generally restrict telling
a customer what screening or reporting has happened about them, and rather than
work out exactly where the line falls in each jurisdiction, the safe operating
rule is that screening detail does not go to its subject --- then the legal
question never has to be reached. Second, it teaches evasion: ``matched at 86\%,
downgraded because the listed date of birth conflicts'' tells someone precisely
which field to change. Third, and most often overlooked, \emph{most of these
matches are wrong}: a name similarity is a finding about a name, not about a
person, and showing it attaches a sanctioned individual's identity to an
applicant who is, in the ordinary case, somebody else entirely.

On the law specifically: in the UK the nearest provision is the Proceeds of
Crime Act 2002 s.333A, but it is narrower than the rule I apply --- it concerns
disclosing that a suspicious activity report has been made or that an
investigation is contemplated. I name it as the family of rule this sits in, not
as authority for it.

The implementation is a \textbf{whitelist that fails closed}. Each event type
maps to a sentence written for the applicant; an action with no entry renders
nothing at all. A new event type added later is invisible to applicants until
somebody decides what that audience may be told, which is the safe direction for
that decision to default in.

One subtlety worth knowing, because it is easy to undo: this works because the
page is a React Server Component, so the raw rows never cross to the browser.
Moving that list into a Client Component would serialise them into the payload
again --- leaking without a single line of rendering code changing. A comment
would not survive that refactor, so there is a test that fetches the live page
and asserts that ``OFAC'', ``name similarity'', the score and the raw payload
keys appear nowhere in the response. It has a counterweight test asserting the
page still renders the applicant's own history, because a disclosure assertion
is trivially satisfied by a blank page.

\section{Sanctions provenance: what the hash does and does not prove}

A decision has two inputs: the rules, and the list. I versioned the rules early.
The list carried no version at all, so the question an auditor actually asks ---
\emph{``was this person screened against the list as it stood that day?''} ---
was not hard to answer, it was unanswerable.

Now every screening result points at a snapshot row recording the source, the
publisher's own publication date, the record count, a SHA-256 of the file, and
when it was loaded. The chain runs decision → match → list version → digest.

\textbf{Precisely what the hash proves:} that two screening runs recording the
same digest used byte-identical content. That is useful --- it makes ``the same
list'' checkable rather than assumed.

\textbf{What it does not prove.} The digest is of the JSON file my downloader
writes, not of the XML the US Treasury published, so it attests to my
\emph{parser's output}, not to the publisher's document. The file is gitignored
and never committed, so a reader cannot recompute it from the repository. And it
is recorded, not verified: nothing checks the digest on load, so a corrupted file
would be recorded faithfully as corrupt. The right guard is cross-checking
against OFAC's own declared record count and the previous snapshot, which I have
not built.

Staleness is also unsolved. The list is loaded at deployment and nothing
refreshes it. Recording the publication date makes that \emph{visible} --- an old
date next to a recent decision is readable evidence --- and visible is a real
improvement over invisible, but it is not the same as fixed.

\section{Two-stage matching, and one principle about pre-filters}

Fuzzy name matching against 19{,}393 OFAC entries (44{,}017 searchable spellings
once aliases are included) originally happened in memory: 86\,MB resident,
scanned per applicant. That is fine for a process that runs for weeks and
impossible for one that must start, answer and exit.

So the list moved into Postgres and the search became two stages. A trigram GIN
index over individual name \emph{tokens} narrows the list to a shortlist ---
averaging 27 entries, 0.141\% of the list --- and then rapidfuzz scores that
shortlist with exactly the code that ran before.

I indexed tokens rather than whole names, and that was the real decision.
Measured against pairs the matcher \emph{accepts}, whole-string trigram
similarity falls as low as 0.471 (\texttt{ahmed hassan} against \texttt{ahmad
hasan}), because transliteration drift lives inside words and an extra middle
name wrecks a whole-string comparison while leaving every other token identical.
Per-token similarity for the same pairs stays at 0.556 and above.

\subsection{The principle}

My first version of the shortlist required two trigram-matched tokens, mirroring
the Python rule that two name parts must correspond. It looked right and was
wrong, because it applied that rule using a \emph{different similarity metric}
from the one the scorer uses, and the two disagree badly on short tokens and
substrings: \texttt{chen}/\texttt{jicheng} scores 72.7 on the scorer's metric
(accepted) and 0.22 on trigram similarity (rejected).

Thirteen applicants out of 400 lost a real match. Every one was a miss, never a
false positive --- which is the dangerous direction.

\textbf{A cheap pre-filter must only ever widen.} The moment it reimplements the
expensive filter's judgement with a cheaper metric, it can disagree with it, and
every disagreement is a match nobody sees. The shortlist now asks the weakest
useful question --- does this entry share any plausible name part? --- and leaves
every actual decision to the scorer.

The threshold was then set by measurement rather than taste: 0.35 lost 13
applicants' matches, 0.30 lost one, 0.25 loses none. I chose 0.25 because recall
is worth more than shortlist size here --- a larger shortlist costs microseconds
of Python, and a shortlist missing the one true match costs a sanctioned person
an account.

\section{Differential testing, and the bug that was already shipping}

Changing how candidates are \emph{found} must not change who gets flagged. A
silent shift in that is the worst outcome available: nobody sees it, the tests
stay green, and the only evidence is a different set of people being refused.

So before trusting the new matcher I wrote an equivalence test that runs both
over 400 seeded applicants and compares the matched entities, the strengths, the
matched spelling, the date-of-birth conflict flag, the resulting score and the
routing. It found three real bugs, all of which would have shipped.

One of them was not mine. An applicant matched three OFAC entities at
\emph{exactly} 100.00 --- two with no date-of-birth conflict, one with. The
scorer takes the strongest hit of each type, which with \texttt{max()} means the
first of several equals. So whichever entity happened to be listed first decided
whether the match was downgraded: the same applicant, against the same list,
scored 60 or 75 depending on dictionary iteration order.

That was already in production behaviour before the migration. Running two
implementations side by side is what made it visible; neither implementation was
wrong on its own, and no single-implementation test would have caught it,
because there was nothing to compare against. Both now sort by score and then by
entity id, which makes the tie-break a decision rather than an accident.

I include this because ``why should I trust your judgement'' is a fair question
and the honest answer is not that I do not make mistakes. It is that I try to
build the thing that would tell me.

\section{When one service may call another}

The two services communicate only through the database. That rule exists to
avoid an internal API --- one service asking the other to do something and
waiting to hear how it went, which drags CORS, shared auth, retries and a second
failure mode into a design that had none.

The deployed version bends it: after enqueueing work, the web application sends
an HTTP request to the worker so an applicant sees a decision in seconds rather
than waiting for the next scheduled run. I let myself do that only because it
passes a specific test:

\begin{itemize}\itemsep0em
\item the work is already durable --- committed to the queue --- before the call
\item the call carries no payload and expects no answer
\item it is not awaited, and cannot fail the request that triggered it
\item a scheduled path does the same work if it never arrives
\end{itemize}

Which reduces to one sentence: \textbf{delete the call and the system is still
correct, only slower}. The queue still carries all the meaning; the HTTP request
carries only timing.

That test then proved itself in an unplanned way. The call silently did nothing
for a period after deployment --- and nothing broke. Applications were still
decided, by the scheduled path, a few minutes later instead of a few seconds.
The failure was a latency regression rather than a correctness one, which is
exactly what the design predicted.

\section{The worker as an on-demand function}

Locally the worker is a long-lived loop: claim, run, sleep, repeat. That is the
right shape and it is what \texttt{docker compose up} gives you. It is also a
process that must run somewhere permanently, and permanently costs money. I am
not paying a monthly bill to host a portfolio demo, and I think that is a
sufficient reason.

So production has a second entry point: an HTTP endpoint that claims a
\emph{bounded} batch --- five jobs --- runs them and returns. Bounded because a
function has a wall-clock limit; a loop that ran until the queue emptied would be
killed mid-job on every invocation once a backlog existed, making no progress
while looking busy. It is called repeatedly instead: once by the web app when
work is enqueued, and every five minutes by a scheduler, which is also what runs
the two recurring jobs a long-lived process would have run on a timer.

The part I care about is that the queue's guarantees are untouched. Both entry
points import the same function to execute a job, so the transaction boundary and
the retry accounting cannot diverge. Jobs are still claimed one at a time with
the same statement, and the 500-job no-double-processing proof runs against both.
Same logic, two runtimes --- which is the honest description, rather than
pretending the deployed thing is the same thing.

\textbf{The cost:} an applicant arriving when the database has been idle waits
about 2.4 seconds for the first page that queries it, and roughly 0.3 seconds for
every page after. In practice the five-minute schedule keeps the database awake,
so this is only paid if the schedule has stopped.

\vfill
\rule{\textwidth}{0.4pt}
{\footnotesize Synthetic data throughout. No real person's data has been in this
system. Figures in this document were measured, not estimated; where a number
moves between runs I have said so.}

\end{document}
